\slcol{श्रीभगवानुवाच —\\
\Index{अनाश्रितः कर्मफलं} कार्यं कर्म करोति यः ।\\
स संन्यासी च योगी च न निरग्निर्न चाक्रियः ॥ ६.१ ॥}
\translate{śrībhagavānuvāca —\\
anāśritaḥ karmaphalaṁ kāryaṁ karma karōti yaḥ ।\\
sa saṁnyāsī ca yōgī ca na niragnirna cākriyaḥ ॥ 6.1 ॥}
\cquote{Sri Bhagavan said:\\
One who performs his duty without depending on the fruits of his actions, is a true sannyasi and a yogi, not the one who forsakes fire ritual or ceases all activity.} 

\slcol{\Index{यं संन्यासमिति} प्राहुर्योगं तं विद्धि पाण्डव ।\\
न ह्यसंन्यस्तसङ्कल्पो योगी भवति कश्चन ॥ ६.२ ॥}
\translate{yaṁ saṁnyāsamiti prāhuryōgaṁ taṁ viddhi pāṇḍava ।\\
na hyasaṁnyastasaṅkalpō yōgī bhavati kaścana ॥ 6.2 ॥}
\cquote{Arjuna, Renunciation is indeed what is known as yoga. No one can become a yogi without renouncing all desires.}


\newpage
\begin{mananam}{\mananamfont {Mananam sloka - 1, 2}}
\mananamtext When I consider giving up something or resigning from a role, what drives my motivation? Is it a desire for an easier life without discomfort? Alternatively, am I motivated by the prospect of seeking better rewards? Can I find motivation within myself to make efforts purely out of duty or a desire to provide unconditional help? Is it possible for me to cultivate a higher trust that rewards will come when I truly deserve them?
\end{mananam}
\WritingHand\enspace{\selfReflect\small{Self Reflection}}
\begin{inspiration}{\mananamfont Inspiration}
\mananamtext The ancient concept of a sannyasi is someone who has renounced all Vedic rituals. Why do they renounce these rituals? Because they no longer seek any material or higher worldly benefits. The Lord clarifies that the true essence of sannyasa, the highest calling in life, is to renounce all personal expectations from our efforts, not the actions themselves. Actions, in themselves, are not a hindrance to spiritual progress; it is the focus on the results that are the obstacles.
\end{inspiration}
\newpage

\slcol{\Index{आरुरुक्षोर्मुनेर्योगं} कर्म कारणमुच्यते ।\\
योगारूढस्य तस्यैव शमः कारणमुच्यते ॥ ६.३ ॥}
\translate{yōgārūḍhasya tasyaiva śamaḥ kāraṇamucyatē ।\\
ārurukṣōrmunēryōgaṁ karma kāraṇamucyatē ॥ 6.3 ॥}
\cquote{For the sage who seeks to ascend to yoga, action is considered to be the means. For that person, who has ascended to yoga, renunciation of action alone is said to be the means.}

\slcol{\Index{यदा हि नेन्द्रियार्थेषु} न कर्मस्वनुषज्जते ।\\
सर्वसङ्कल्पसंन्यासी योगारूढस्तदोच्यते ॥ ६.४ ॥ }
\translate{yadā hi nēndriyārthēṣu na karmasvanuṣajjatē ।\\
sarvasaṅkalpasaṁnyāsī yōgārūḍhastadōcyatē ॥ 6.4 ॥}
\cquote{When a man is not attached to sense objects or to the actions, having renounced all desires and intentions, then he is said to have attained yoga.}

\slcol{\Index{उद्धरेदात्मनात्मानं} नात्मानमवसादयेत् ।\\
आत्मैव ह्यात्मनो बन्धुरात्मैव रिपुरात्मनः ॥ ६.५ ॥} 
\translate{uddharēdātmanātmānaṁ nātmānamavasādayēt ।\\
ātmaiva hyātmanō bandhurātmaiva ripurātmanaḥ ॥ 6.5 ॥}
\cquote{One should elevate oneself by one’s own Self and should not degrade oneself. For the self is one’s own friend, and the self is also one’s own foe.}

\slcol{\Index{बन्धुरात्मात्मनस्तस्य} येनात्मैवात्मना जितः ।\\
अनात्मनस्तु शत्रुत्वे वर्तेतात्मैव शत्रुवत् ॥ ६.६ ॥}
\translate{bandhurātmātmanastasya yēnātmaivātmanā jitaḥ ।\\
anātmanastu śatrutvē vartētātmaiva śatruvat ॥ 6.6 ॥}
\cquote{The self is the friend to one who has gained control over it. But for one who lacks self-mastery, the self is an enemy.}



\newpage
\begin{mananam}{\mananamfont {Mananam sloka - 3, 4}}
\mananamtext Is my mind cooperative and willing to sit for meditation? How much time can I regularly dedicate to meditation without feeling lost or overly restless? If I am unable to meditate, can I use that time for selfless activities? How can I balance my daily activities with meditation? How do I balance my activities with meditation? How do I keep my activities in check, to ensure that I don’t compromise my meditation time?
\end{mananam}
\WritingHand\enspace{\selfReflect\small{Self Reflection}}
\begin{inspiration}{\mananamfont Inspiration}
\mananamtext A true spiritual teacher or teaching provides personalized instructions to students based on their stage in the spiritual journey. When one has calmed the mind and is free from restlessness and worldly thoughts, that person should cease all physical and mental efforts to enjoy the stillness of meditation. However, if one is not yet ready to sit still for long periods of meditation, that person should engage in selfless actions.
\end{inspiration}
\newpage


\newpage
\begin{mananam}{\mananamfont {Mananam sloka - 5, 6}}
\mananamtext Am I attuned to my conscience, my true friend and guide, in my daily life? Do I adhere to my positive intentions and resolutions, or do I oppose this inner voice of guidance?\\
Do I rebel against well-meaning people in my life, akin to a rebellious teenager? Can I recognize that when I resist the guidance of teachers and the Scriptures, this resistance impedes my own spiritual progress? Am I aware of the ego’s force at play in my life?
\end{mananam}
\WritingHand\enspace{\selfReflect\small{Self Reflection}}
\begin{inspiration}{\mananamfont Inspiration}
\mananamtext When we attune to our true Self, which embodies pure Consciousness and Bliss, our mind becomes positive and uplifted. Such a mind collaborates with our intelligence, fostering overall well-being. However, when the mind is swayed by the senses, it pursues transient worldly pleasures and gratification. In doing so, it rebels against the intelligence attuned to Consciousness, adhering to the adage ‘Ignorance is bliss.’ Unfortunately, such a life path is fraught with suffering.
\end{inspiration}
\newpage

\slcol{\Index{जितात्मनः प्रशान्तस्य} परमात्मा समाहितः ।\\
शीतोष्णसुखदुःखेषु तथा मानापमानयोः ॥ ६.७ ॥}
\translate{jitātmanaḥ praśāntasya paramātmā samāhitaḥ ।\\
śītōṣṇasukhaduḥkhēṣu tathā mānāpamānayōḥ ॥ 6.7 ॥}
\cquote{The supreme self, in one who is self-controlled and peaceful, remains balanced in cold and heat, in pleasure and pain, as well as in honour and dishonour.}
\slcol{\Index{ज्ञानविज्ञानतृप्तात्मा} कूटस्थो विजितेन्द्रियः ।\\
युक्त इत्युच्यते योगी समलोष्टाश्मकाञ्चनः ॥ ६.८ ॥}
\translate{jñānavijñānatr̥ptātmā kūṭasthō vijitēndriyaḥ ।\\
yukta ityucyatē yōgī samalōṣṭāśmakāñcanaḥ ॥ 6.8 ॥}
\cquote{He who is satisfied with both knowledge and experience, who is steady and has mastered the senses, is united in yoga. Such a person, who sees a lump of clay, a stone and gold equally is a yogi.}


\slcol{\Index{सुहृन्मित्रार्युदासीनम}ध्यस्थद्वेष्यबन्धुषु ।\\
साधुष्वपि च पापेषु समबुद्धिर्विशिष्यते ॥ ६.९ ॥}
\translate{suhr̥nmitrāryudāsīnamadhyasthadvēṣyabandhuṣu ।\\
sādhuṣvapi ca pāpēṣu samabuddhirviśiṣyatē ॥ 6.9 ॥}
\cquote{He who looks with equal regard upon a friend, a companion, a noble person, a detached, the hostile, and a relative - and who is impartial towards both the righteous and the unrighteous attains excellence.}

\slcol{\Index{योगी युञ्जीत} सततमात्मानं रहसि स्थितः ।\\
एकाकी यतचित्तात्मा निराशीरपरिग्रहः ॥ ६.१० ॥}
\translate{yōgī yuñjīta satatamātmānaṁ rahasi sthitaḥ ।\\
ēkākī yatacittātmā nirāśīraparigrahaḥ ॥ 6.10 ॥}
\cquote{A yogi should constantly concentrate through meditation, remaining in solitude alone, with a focused mind and self-controlled mind, free from desire, expectation and possessions.}



\newpage
\begin{mananam}{\mananamfont {Mananam sloka - 8, 9}}
\mananamtext Do I perceive the one Divine essence in everything, or do I view things and people solely through the lens of my mental inclinations and past experiences? To what extent do my memories influence and colour my interactions with others? What benefits lie in adopting a fresh perspective toward every experience in my life? What obstacles prevent me from putting that fresh perspective into action?
\end{mananam}
\WritingHand\enspace{\selfReflect\small{Self Reflection}}
\begin{inspiration}{\mananamfont Inspiration}
\mananamtext Realized yogis perceive the underlying Consciousness in everything. To them, everyone appears as a spark of this One Divine Consciousness. Consequently, their interactions are free from prejudices and biases. In contrast, for everyone else, the world is marked by labels, names, categories, and defined shapes that distinguish one object from another.
\end{inspiration}
\newpage


\newpage
\begin{mananam}{\mananamfont {Mananam sloka - 10}}
\mananamtext Do I allocate time in my daily life to simply be by myself? Am I comfortable with occasional solitude, or does it make me uneasy? Do I experience anxiety, fear, boredom, or a sense of missing out during these moments?\\
When alone, do I feel compelled to use music, TV, or other gadgets as substitutes for human interaction? What spiritual and mental benefits can an hour of solitude and contemplation offer me?
\end{mananam}
\WritingHand\enspace{\selfReflect\small{Self Reflection}}
\begin{inspiration}{\mananamfont Inspiration}
\mananamtext The modern world differs significantly from the lifestyle depicted in ancient India. Technology and social media have deprived humans of the opportunity for solitude. Solitude allows for self-reflection, inner peace, and mental rejuvenation; it fosters creativity, clarity of thought, and a deeper understanding of oneself.\\
Learning to be with oneself in a positive manner is a valuable skill—an art worth mastering. Consistently practicing this skill is akin to a science. The ancient rishis named this art and science, ‘Meditation.’
\end{inspiration}
\newpage

\slcol{\Index{शुचौ देशे प्रतिष्ठाप्य} स्थिरमासनमात्मनः ।\\
नात्युच्छ्रितं नातिनीचं चैलाजिनकुशोत्तरम् ॥ ६.११ ॥}
\translate{śucau dēśē pratiṣṭhāpya sthiramāsanamātmanaḥ ।\\
nātyucchritaṁ nātinīcaṁ cailājinakuśōttaram ॥ 6.11 ॥}
\cquote{Having established a clean and firm seat, neither too high nor too low, made of cloth, deerskin and kush grass placed one over the other in succession.}
\slcol{\Index{तत्रैकाग्रं मनः} कृत्वा यतचित्तेन्द्रियक्रियः ।\\
उपविश्यासने युञ्ज्याद्योगमात्मविशुद्धये ॥ ६.१२ ॥}
\translate{tatraikāgraṁ manaḥ kr̥tvā yatacittēndriyakriyaḥ ।\\
upaviśyāsanē yuñjyādyōgamātmaviśuddhayē ॥ 6.12 ॥}
\cquote{There, sitting on that seat, with the mind concentrated and the mind and senses held in check, one should practice yoga for the purification of the self.}

\slcol{\Index{समं कायशिरोग्रीवं} धारयन्नचलं स्थिरः ।\\
सम्प्रेक्ष्य नासिकाग्रं स्वं दिशश्चानवलोकयन् ॥ ६.१३ ॥} 
\translate{samaṁ kāyaśirōgrīvaṁ dhārayannacalaṁ sthiraḥ ।\\
samprēkṣya nāsikāgraṁ svaṁ diśaścānavalōkayan ॥ 6.13 ॥}
\cquote{Holding firmly body, head, neck erect and still, with the eye-balls fixed, as if gazing at the tip of the nose, and without looking around.}

\slcol{\Index{प्रशान्तात्मा विगतभी}र्ब्रह्मचारिव्रते स्थितः ।\\
मनः संयम्य मच्चित्तो युक्त आसीत मत्परः ॥ ६.१४ ॥}
\translate{praśāntātmā vigatabhīrbrahmacārivratē sthitaḥ ।\\
manaḥ saṁyamya maccittō yukta āsīta matparaḥ ॥ 6.14 ॥}
\cquote{Serene and fearless, firmly established  in the vow of celibacy, with the mind controlled and fixed on Me, the yogi should sit, having Me as the supreme goal.}



\newpage
\begin{mananam}{\mananamfont {Mananam sloka - 12}}
\mananamtext
Suggested areas of Meditation Introspection for those aspiring for a regular practice are:\\
\textbf{Consistent Practice}: Do I maintain a regular meditation practice in my life?\\
\textbf{Sensory Restraint}: Can I keep my eyes closed, other senses withdrawn, and maintain stillness in my body all through the meditation period?\\
\textbf{Focus}: How quickly am I able to achieve a state of one-pointedness?\\
\textbf{State of Mind}: Is my mind often restless or dull?
Process: Do I comprehend the meditation process well? Am I open to seeking guidance from a teacher?
\end{mananam}
\WritingHand\enspace{\selfReflect\small{Self Reflection}}
\begin{inspiration}{\mananamfont Inspiration}
\mananamtext Meditation is akin to inner purification of the mind—a vital practice, much like bathing to keep the body clean. Mastering the formal meditation process is essential to unlock the benefits of peace and inner happiness.\\
According to the Bhagavad Gita, an ideal yogi isn’t someone who bends nature to his will or performs miracles. Instead, the yogi learns to master his own responses to the world by controlling the senses and mind.
\end{inspiration}
\newpage

\slcol{\Index{युञ्जन्नेवं सदात्मानं} योगी नियतमानसः ।\\
शान्तिं निर्वाणपरमां मत्संस्थामधिगच्छति ॥ ६.१५ ॥}
\translate{yuñjannēvaṁ sadātmānaṁ yōgī niyatamānasaḥ ।\\
śāntiṁ nirvāṇaparamāṁ matsaṁsthāmadhigacchati ॥ 6.15 ॥}
\cquote{Thus, constantly engaging the self in meditation, and with a controlled mind, the yogi attains peace abiding in Me, which leads to liberation (Nirvana).}
\slcol{\Index{नात्यश्नतस्तु योगोऽस्ति} न चैकान्तमनश्नतः ।\\
न चातिस्वप्नशीलस्य जाग्रतो नैव चार्जुन ॥ ६.१६ ॥}
\translate{nātyaśnatastu yōgō:'sti na caikāntamanaśnataḥ ।\\
na cātisvapnaśīlasya jāgratō naiva cārjuna ॥ 6.16 ॥}
\cquote{O Arjuna, yoga is not attainable by one who overeats nor by one who fasts too excessively; neither by one who sleeps too much nor by one who stays awake excessively.}


\slcol{\Index{युक्ताहारविहारस्य} युक्तचेष्टस्य कर्मसु ।\\
युक्तस्वप्नावबोधस्य योगो भवति दुःखहा ॥ ६.१७ ॥}
\translate{yuktāhāravihārasya yuktacēṣṭasya karmasu ।\\
yuktasvapnāvabōdhasya yōgō bhavati duḥkhahā ॥ 6.17 ॥}
\cquote{For one who regulates their food and recreation, is balanced in actions, and maintains moderation in sleep and wakefulness, Yoga will dispel all unhappiness.}

\slcol{\Index{यदा विनियतं} चित्तमात्मन्येवावतिष्ठते ।\\
निःस्पृहः सर्वकामेभ्यो युक्त इत्युच्यते तदा ॥ ६.१८ ॥}
\translate{yadā viniyataṁ cittamātmanyēvāvatiṣṭhatē ।\\
niḥspr̥haḥ sarvakāmēbhyō yukta ityucyatē tadā ॥ 6.18 ॥}
\cquote{When one’s controlled mind rests in the Self alone, free from all desires for worldly objects, that person is considered truly united in yoga.}



\newpage
\begin{mananam}{\mananamfont {Mananam sloka - 17}}
\mananamtext 
Suggested areas of Introspection for achieving moderation in daily life are:\\
\textbf{Diet Balance}: What constitutes a balanced diet for me? Do I sometimes binge on food and then compensate with excessive fasting?\\
\textbf{Activity Levels}: Am I balanced in terms of physical activity? Do I oscillate between being lazy and procrastinating, and then becoming too restless and overactive?\\
\textbf{Sleep Patterns}: How consistent are my sleep patterns? Do I sometimes oversleep and at other times get insufficient rest?
\end{mananam}
\WritingHand\enspace{\selfReflect\small{Self Reflection}}
\begin{inspiration}{\mananamfont Inspiration}
\mananamtext Long-term success on the spiritual path lies in maintaining a balanced life. Practitioners must be mindful not to overindulge their senses, but equally important is avoiding excessive deprivation through food and sleep, which weakens the body. A healthy body serves as a crucial vehicle for spiritual pursuits. The Bhagavad Gita advocates embracing the middle path and avoiding extremes.
\end{inspiration}
\newpage

\slcol{\Index{यदा दीपो निवातस्थो} नेङ्गते सोपमा स्मृता ।\\
योगिनो यतचित्तस्य युञ्जतो योगमात्मनः ॥ ६.१९ ॥}
\translate{yadā dīpō nivātasthō nēṅgatē sōpamā smr̥tā ।\\
yōginō yatacittasya yuñjatō yōgamātmanaḥ ॥ 6.19 ॥}
\cquote{Just as a lamp placed in a windless spot does not flicker, so is the yogi compared - one whose mind is disciplined and who is fully engaged in yoga for self-realization.}
\slcol{\Index{यत्रोपरमते चित्तं} निरुद्धं योगसेवया ।\\
यत्र चैवात्मनात्मानं पश्यन्नात्मनि तुष्यति ॥ ६.२० ॥} 
\translate{yatrōparamatē cittaṁ niruddhaṁ yōgasēvayā ।\\
yatra caivātmanātmānaṁ paśyannātmani tuṣyati ॥ 6.20 ॥}
\cquote{When, through the discipline of yoga, the mind becomes still, and one perceives the Self with the Self, then one feels contentment within the Self.}

\slcol{\Index{सुखमात्यन्तिकं} यत्तद्बुद्धिग्राह्यमतीन्द्रियम् ।\\
वेत्ति यत्र न चैवायं स्थितश्चलति तत्त्वतः ॥ ६.२१ ॥}
\translate{sukhamātyantikaṁ yattadbuddhigrāhyamatīndriyam ।\\
vētti yatra na caivāyaṁ sthitaścalati tattvataḥ ॥ 6.21 ॥}
\cquote{Where one realizes the boundless bliss that is understood through the intellect and lies beyond the senses, and once established in it, one does not deviate from that truth.}


\slcol{\Index{यं लब्ध्वा चापरं} लाभं मन्यते नाधिकं ततः ।\\
यस्मिन्स्थितो न दुःखेन गुरुणापि विचाल्यते ॥ ६.२२ ॥}
\translate{yaṁ labdhvā cāparaṁ lābhaṁ manyatē nādhikaṁ tataḥ ।\\
yasminsthitō na duḥkhēna guruṇāpi vicālyatē ॥ 6.22 ॥}
\cquote{Upon attaining that state, one does not consider any gain superior, and once established in it, even the greatest misery cannot disturb their inner stability.}

\newpage
\slcol{\Index{तं विद्याद्दुःखसंयोग}वियोगं योगसंज्ञितम् ।\\
स निश्चयेन योक्तव्यो योगोऽनिर्विण्णचेतसा ॥ ६.२३ ॥} 
\translate{taṁ vidyādduḥkhasaṁyōgaviyōgaṁ yōgasaṁjñitam ।\\
sa niścayēna yōktavyō yōgō:'nirviṇṇacētasā ॥ 6.23 ॥}
\cquote{This practice, which involves breaking the bond with suffering, is called yoga. It should be pursued with firm resolve and a mind free from discouragement.}
\slcol{\Index{सङ्कल्पप्रभवान्कामां}स्त्यक्त्वा सर्वानशेषतः ।\\
मनसैवेन्द्रियग्रामं विनियम्य समन्ततः ॥ ६.२४ ॥}
\translate{saṅkalpaprabhavānkāmāṁstyaktvā sarvānaśēṣataḥ ।\\
manasaivēndriyagrāmaṁ viniyamya samantataḥ ॥ 6.24 ॥}
\cquote{Having renounced all desires born from intentions, and fully restraining the group of sense organs with the mind from all sides.}

\slcol{\Index{शनैः शनैरुपरमेद्बुद्ध्या} धृतिगृहीतया ।\\ 
आत्मसंस्थं मनः कृत्वा न किञ्चिदपि चिन्तयेत् ॥ ६.२५ ॥} 
\translate{śanaiḥ śanairuparamēdbuddhyā dhr̥tigr̥hītayā ।\\
ātmasaṁsthaṁ manaḥ kr̥tvā na kiñcidapi cintayēt ॥ 6.25 ॥}
\cquote{Gradually, one attains calmness through an intellect supported by unwavering determination. Once the mind is established in the Self, nothing else is to be considered.}

\slcol{\Index{यतो यतो निश्चरति} मनश्चञ्चलमस्थिरम् ।\\ 
ततस्ततो नियम्यैतदात्मन्येव वशं नयेत् ॥ ६.२६ ॥ }
\translate{yatō yatō niścarati manaścañcalamasthiram ।\\
tatastatō niyamyaitadātmanyēva vaśaṁ nayēt ॥ 6.26 ॥}
\cquote{Whenever the restless and unsteady mind wanders, one should subdue it immediately and bring it back under the control of the Self.}



\newpage
\begin{mananam}{\mananamfont {Mananam sloka - 22, 23}}
\mananamtext Why do I engage in yoga and other well-being practices? Is it for earning a livelihood, social status, recognition, pride, etc.? Or is it for being free of pain and sorrows – physically, mentally, emotionally, etc.? Am I doing it for my own well-being or to impress others? Am I convinced that this path leads to the highest level of well-being?
\end{mananam}
\WritingHand\enspace{\selfReflect\small{Self Reflection}}
\begin{inspiration}{\mananamfont Inspiration}
\mananamtext The cessation of all mental suffering and the attainment of sublime happiness is the goal of all spiritual practices and study. Without a clear determination of purpose, practicing and studying these teachings for worldly gains will eventually lead to dissatisfaction and disappointment. Suffering will come to an end when one has this clarity, along with the determination in undertaking the practice. This is the promise of yoga.
\end{inspiration}
\newpage



\newpage
\begin{mananam}{\mananamfont {Mananam sloka - 25, 26}}
\mananamtext Have I clearly determined the object of focus in my meditation practice?\\
Am I consistently maintaining awareness of the mind even when distractions arise? Do I then gently remind the mind to return to object of meditation?\\
Do I remain patient when my mind becomes restless during meditation, or do I give up prematurely?\\
Do I observe the immediate benefits of mind training—such as achieving a quieter state of mind—right after meditation? Am I also able to experience these benefits manifesting later in the day? 
\end{mananam}
\WritingHand\enspace{\selfReflect\small{Self Reflection}}
\begin{inspiration}{\mananamfont Inspiration}
\mananamtext All meditation techniques boil down to a simple formula: Whenever the mind gets distracted, we gently remind ourselves to refocus our attention. This practice combines ancient wisdom with the modern science of attention regulation. Patient perseverance is the essence of this approach, leading to certain success.
\end{inspiration}
\newpage

\slcol{\Index{प्रशान्तमनसं ह्येनं} योगिनं सुखमुत्तमम् ।\\
उपैति शान्तरजसं ब्रह्मभूतमकल्मषम् ॥ ६.२७ ॥}
\translate{praśāntamanasaṁ hyēnaṁ yōginaṁ sukhamuttamam ।\\
upaiti śāntarajasaṁ brahmabhūtamakalmaṣam ॥ 6.27 ॥}
\cquote{Supreme bliss truly comes to the yogi, whose mind is peaceful, whose passions are subdued, who is free from sin and who has become one with the Brahman.}
\slcol{\Index{युञ्जन्नेवं सदात्मानं} योगी विगतकल्मषः । \\
सुखेन ब्रह्मसंस्पर्शमत्यन्तं सुखमश्नुते ॥ ६.२८ ॥ }
\translate{yuñjannēvaṁ sadātmānaṁ yōgī vigatakalmaṣaḥ ।\\
sukhēna brahmasaṁsparśamatyantaṁ sukhamaśnutē ॥ 6.28 ॥ }
\cquote{By constantly concentrating the mind in self-discipline, the sinless yogi attains the infinite bliss of union with Brahman.}


\slcol{\Index{सर्वभूतस्थमात्मानं} सर्वभूतानि चात्मनि ।\\
ईक्षते योगयुक्तात्मा सर्वत्र समदर्शनः ॥ ६.२९ ॥} 
\translate{sarvabhūtasthamātmānaṁ sarvabhūtāni cātmani ।\\
īkṣatē yōgayuktātmā sarvatra samadarśanaḥ ॥ 6.29 ॥}
\cquote{With the mind self-absorbed in yoga, one sees everything equally, perceiving the Self in all beings and all beings in the Self.}

\slcol{\Index{यो मां पश्यति सर्वत्र} सर्वं च मयि पश्यति ।\\
तस्याहं न प्रणश्यामि स च मे न प्रणश्यति ॥ ६.३० ॥}
\translate{yō māṁ paśyati sarvatra sarvaṁ ca mayi paśyati ।\\
tasyāhaṁ na praṇaśyāmi sa ca mē na praṇaśyati ॥ 6.30 ॥}
\cquote{One who sees Me in everything and sees everything in Me, is never separated from Me, and I am never separated from him.}

\slcol{\Index{सर्वभूतस्थितं यो मां} भजत्येकत्वमास्थितः ।\\
सर्वथा वर्तमानोऽपि स योगी मयि वर्तते ॥ ६.३१ ॥ }
\translate{sarvabhūtasthitaṁ yō māṁ bhajatyēkatvamāsthitaḥ ।\\
sarvathā vartamānō:'pi sa yōgī mayi vartatē ॥ 6.31 ॥ }
\cquote{The yogi who sees unity and worships Me as dwelling in all beings remains ever in Me, whatever may be his mode of living.}



\newpage
\begin{mananam}{\mananamfont {Mananam sloka - 30}}
\mananamtext What does God mean to me in my daily life? Can I invoke the presence of God in all circumstances?\\
Do I experience a sense of disconnection from God when I am not in the presence of an image or any tangible reminder of God’s existence?\\
If not, how can I adapt the concept of God to something that provides constant refuge and connection?
\end{mananam}
\WritingHand\enspace{\selfReflect\small{Self Reflection}}
\begin{inspiration}{\mananamfont Inspiration}
\mananamtext Experiencing God’s presence consistently uplifts the mind. It transforms the material world into something Divine, dispelling feelings of being lost and helpless. For a devoted soul, nothing is more comforting than feeling God’s presence in every aspect of life.
\end{inspiration}
\newpage

\newpage
\begin{mananam}{\mananamfont {Mananam sloka - 31, 32}}
\mananamtext Can I recognize and feel my own Self in everyone around me?\\
Beyond transactional purposes, how do I value people—based on their roles and achievements, or rather, upon their inherent God-essence?\\
Can I relate to my own Self beyond my outer roles, skills, and personality?\\
And can this awareness of God-essence in myself and everyone transform how I interact with others in my daily life?
\end{mananam}
\WritingHand\enspace{\selfReflect\small{Self Reflection}}
\begin{inspiration}{\mananamfont Inspiration}
\mananamtext For those on the path of wisdom, a valuable practice is to recognize that the essence of our Divine Consciousness—the Self—is also the essence of all beings. With this understanding, there are no enemies, and fear dissipates, for the entire world is nothing but an extension of our own Being.
\end{inspiration}
\newpage

\slcol{\Index{आत्मौपम्येन सर्वत्र} समं पश्यति योऽर्जुन ।\\
शुखं वा यदि वा दुःखं स योगी परमो मतः ॥ ६.३२ ॥ }
\translate{ātmaupamyēna sarvatra samaṁ paśyati yō:'rjuna ।\\
śukhaṁ vā yadi vā duḥkhaṁ sa yōgī paramō mataḥ ॥ 6.32 ॥}
\cquote{O Arjuna, that yogi is regarded as the highest who views the happiness and sorrow of all beings with the same standard he would apply to himself.}
\slcol{अर्जुन उवाच —\\
\Index{योऽयं योगस्त्वया} प्रोक्तः साम्येन मधुसूदन । \\
एतस्याहं न पश्यामि चञ्चलत्वात्स्थितिं स्थिराम् ॥ ६.३३ ॥ }
\translate{arjuna uvāca —\\
yō:'yaṁ yōgastvayā prōktaḥ sāmyēna madhusūdana । \\
ētasyāhaṁ na paśyāmi cañcalatvātsthitiṁ sthirām ॥ 6.33 ॥ }
\cquote{Arjuna said,\\
O Madhusudana, I do not see how I can attain the equanimity in yoga that you have described, due to the restlessness of the mind.}

\slcol{\Index{चञ्चलं हि मनः} कृष्ण प्रमाथि बलवद्दृढम् ।\\
तस्याहं निग्रहं मन्ये वायोरिव सुदुष्करम् ॥ ६.३४ ॥ }
\translate{cañcalaṁ hi manaḥ kr̥ṣṇa pramāthi balavaddr̥ḍham ।\\
tasyāhaṁ nigrahaṁ manyē vāyōriva suduṣkaram ॥ 6.34 ॥ }
\cquote{The mind is restless, turbulent, strong, and unyielding, O Krishna, I consider controlling it is as difficult as controlling the wind.}

\slcol{श्रीभगवानुवाच —\\
\Index{असंशयं महाबाहो} मनो दुर्निग्रहं चलम् ।\\
अभ्यासेन तु कौन्तेय वैराग्येण च गृह्यते ॥ ६.३५ ॥}
\translate{śrībhagavānuvāca —\\
asaṁśayaṁ mahābāhō manō durnigrahaṁ calam । \\
abhyāsēna tu kauntēya vairāgyēṇa ca gr̥hyatē ॥ 6.35 ॥}
\cquote{Sri Bhagawan said,\\
Undoubtedly, O Mahabaho (mighty-armed Arjuna), the mind is restless and difficult to control, but through practice and detachment, O son of Kunti, it can be mastered.}



\newpage
\begin{mananam}{\mananamfont {Mananam sloka - 34, 35}}
\mananamtext Am I mindful of the restless nature of my mind?\\
Do I pause mentally to acknowledge its impulses and projections instead of being carried away by them?\\
Have I consciously made it a practice to pause during the day?\\
How effective is my regular practice in calming the mind?\\
What is my understanding of dispassion, and how can I apply it to avoid distractions and enhance focus in specific activities?
\end{mananam}
\WritingHand\enspace{\selfReflect\small{Self Reflection}}
\begin{inspiration}{\mananamfont Inspiration}
\mananamtext The mind’s inherent nature is restless, and when we consciously acknowledge this—especially during its turbulent moments—we distance ourselves from it. Dispassion and detachment serve as valuable tools for anyone seeking to control the mind. Even a student aiming to focus on studies must temporarily set aside other amusements and distractions. Through diligent practice, one becomes a yogi—a master of one’s own mind, and the world.
\end{inspiration}
\newpage

\slcol{\Index{असंयतात्मना योगो} दुष्प्राप इति मे मतिः ।\\
वश्यात्मना तु यतता शक्योऽवाप्तुमुपायतः ॥ ६.३६ ॥}
\translate{asaṁyatātmanā yōgō duṣprāpa iti mē matiḥ ।\\
vaśyātmanā tu yatatā śakyō:'vāptumupāyataḥ ॥ 6.36 ॥}
\cquote{Yoga is difficult to attain for one whose mind is not under control; however, in My view, it can be achieved by one who is self-disciplined and strives through the proper means.}
\slcol{अर्जुन उवाच —\\
\Index{अयतिः श्रद्धयोपेतो} योगाच्चलितमानसः ।\\
अप्राप्य योगसंसिद्धिं कां गतिं कृष्ण गच्छति ॥ ६.३७ ॥}
\translate{arjuna uvāca —\\
ayatiḥ śraddhayōpētō yōgāccalitamānasaḥ ।\\
aprāpya yōgasaṁsiddhiṁ kāṁ gatiṁ kr̥ṣṇa gacchati ॥ 6.37 ॥}
\cquote{Arjuna said,\\
O Krishna, what is the fate of one who has faith but lacks self-discipline, whose mind wanders from yoga and fails to reach perfection?}

\slcol{\Index{कच्चिन्नोभयविभ्रष्ट}श्छिन्नाभ्रमिव नश्यति ।\\
अप्रतिष्ठो महाबाहो विमूढो ब्रह्मणः पथि ॥ ६.३८ ॥}
\translate{kaccinnōbhayavibhraṣṭaśchinnābhramiva naśyati ।\\
apratiṣṭhō mahābāhō vimūḍhō brahmaṇaḥ pathi ॥ 6.38 ॥}
\cquote{O mighty-armed Krishna, having failed in both paths, and deluded on the path of Brahman, will one be lost, like a cloud torn by the wind, with no support?}

\slcol{\Index{एतन्मे संशयं कृष्ण} च्छेत्तुमर्हस्यशेषतः ।\\
त्वदन्यः संशयस्यास्य च्छेत्ता न ह्युपपद्यते ॥ ६.३९ ॥} 
\translate{ētanmē saṁśayaṁ kr̥ṣṇa cchēttumarhasyaśēṣataḥ ।\\
tvadanyaḥ saṁśayasyāsya cchēttā na hyupapadyatē ॥ 6.39 ॥ }
\cquote{O Krishna, you alone can completely dispel this doubt of mine, for there is no one else who can destroy it.}



\newpage
\begin{mananam}{\mananamfont {Mananam sloka - 37, 38}}
\mananamtext What motivates me to seek control over my mind?\\
Am I conscious of how an agitated and unhappy mind can be the source of all life’s problems?\\
Do I tend to give up prematurely without actively attempting to overcome my weaknesses and unpleasant traits?\\
Do I fear that achieving the desired inner state is an insurmountable task, perhaps requiring more than one lifetime?
\end{mananam}
\WritingHand\enspace{\selfReflect\small{Self Reflection}}
\begin{inspiration}{\mananamfont Inspiration}
\mananamtext The lives of saints and sages, who have overcome adversity and obstacles, serve as inspiration. However, for those starting on the spiritual path, this can feel overwhelming. Lack of self-belief and faith may lead some to abandon their pursuit of the loftiest goals. Yet, having glimpsed the potential of inner freedom, one cannot be content with mere worldly pleasures.
\end{inspiration}
\newpage

\slcol{श्रीभगवानुवाच —\\
\Index{पार्थ नैवेह नामुत्र} विनाशस्तस्य विद्यते ।\\
न हि कल्याणकृत्कश्चिद्दुर्गतिं तात गच्छति ॥ ६.४० ॥ }
\translate{śrībhagavānuvāca —\\
pārtha naivēha nāmutra vināśastasya vidyatē ।\\
na hi kalyāṇakr̥tkaściddurgatiṁ tāta gacchati ॥ 6.40 ॥}
\cquote{Sri Bhagavan said,\\
O Partha, for one who performs good deeds, there is no destruction either in this world or the next. O my son, such a person never meets with an unfortunate end.}

\slcol{\Index{प्राप्य पुण्यकृतां लोका}नुषित्वा शाश्वतीः समाः ।\\
शुचीनां श्रीमतां गेहे योगभ्रष्टोऽभिजायते ॥ ६.४१ ॥ }
\translate{prāpya puṇyakr̥tāṁ lōkānuṣitvā śāśvatīḥ samāḥ ।\\
śucīnāṁ śrīmatāṁ gēhē yōgabhraṣṭō:'bhijāyatē ॥ 6.41 ॥}
\cquote{Having attained the worlds of the righteous and dwelling there for many years, one who has fallen from yoga is reborn in the house of the pure and prosperous.}

\slcol{\Index{अथवा योगिनामेव} कुले भवति धीमताम् ।\\
एतद्धि दुर्लभतरं लोके जन्म यदीदृशम् ॥ ६.४२ ॥} 
\translate{athavā yōgināmēva kulē bhavati dhīmatām ।\\
ētaddhi durlabhataraṁ lōkē janma yadīdr̥śam ॥ 6.42 ॥}
\cquote{Or one is born in a family of enlightened yogis. Such a birth is truly rare to attain in this world.}

\slcol{\Index{तत्र तं बुद्धिसंयोगं} लभते पौर्वदेहिकम् ।\\ 
यतते च ततो भूयः संसिद्धौ कुरुनन्दन ॥ ६.४३ ॥}
\translate{tatra taṁ buddhisaṁyōgaṁ labhatē paurvadēhikam ।\\
yatatē ca tatō bhūyaḥ saṁsiddhau kurunandana ॥ 6.43 ॥}
\cquote{In that birth, he regains the wisdom accumulated in his previous life, and with that, he strives even more sincerely for spiritual perfection, O son of Kuru.}



\newpage
\begin{mananam}{\mananamfont {Mananam sloka - 40}}
\mananamtext Do I tend to assess my achievements or failures using extreme judgments? For instance, do I view things as either complete success or total failure? Are my evaluations based on comparing myself to others, or do I set reasonable goals for myself?\\
Can I identify present-day skills or accomplishments that originated from efforts I put in years ago? Just as past labours contribute to current abilities, which of my present-day efforts could potentially lead to future growth and increased mental skills? Can I reasonably deduce that none of my efforts ever go to waste, including my efforts at meditation and spirituality?
\end{mananam}
\WritingHand\enspace{\selfReflect\small{Self Reflection}}
\begin{inspiration}{\mananamfont Inspiration}
\mananamtext Life cannot be measured solely by the world’s labels of success or failure. Likewise, we should not abandon efforts to improve just because we may not achieve our goals in this lifetime. Every endeavour to change ourselves or serve others holds value and will fructify at an appropriate time. According to the Indian Scriptures, life after life, we continue from where we left off previously. Therefore, until our last breath, let us strive for inner perfection.
\end{inspiration}
\newpage


\slcol{\Index{पूर्वाभ्यासेन तेनैव} ह्रियते ह्यवशोऽपि सः ।\\ 
जिज्ञासुरपि योगस्य शब्दब्रह्मातिवर्तते ॥ ६.४४ ॥}
\translate{pūrvābhyāsēna tēnaiva hriyatē hyavaśō:'pi saḥ ।\\
jijñāsurapi yōgasya śabdabrahmātivartatē ॥ 6.44 ॥}
\cquote{By the very power of previous practice, one is carried forward, even if unwilling. Even a person who simply desires to learn yoga goes beyond those who perform Vedic rituals.}
\slcol{\Index{प्रयत्नाद्यतमानस्तु योगी} संशुद्धकिल्बिषः ।\\
अनेकजन्मसंसिद्धस्ततो याति परां गतिम् ॥ ६.४५ ॥ }
\translate{prayatnādyatamānastu yōgī saṁśuddhakilbiṣaḥ ।\\
anēkajanmasaṁsiddhastatō yāti parāṁ gatim ॥ 6.45 ॥}
\cquote{The yogi, striving with determination and earnest effort, purified of all sins and perfected through many lifetimes, ultimately attains the supreme goal.}

\slcol{\Index{तपस्विभ्योऽधिको योगी} ज्ञानिभ्योऽपि मतोऽधिकः ।\\
कर्मिभ्यश्चाधिको योगी तस्माद्योगी भवार्जुन ॥ ६.४६ ॥ }
\translate{tapasvibhyō:'dhikō yōgī jñānibhyō:'pi matō:'dhikaḥ ।\\
karmibhyaścādhikō yōgī tasmādyōgī bhavārjuna ॥ 6.46 ॥ }
\cquote{The yogi is superior to those who practice austerities, those who possess knowledge and to those who perform actions. Therefore, O Arjuna, be a yogi.}

\slcol{\Index{योगिनामपि सर्वेषां} मद्गतेनान्तरात्मना ।\\
श्रद्धावान्भजते यो मां स मे युक्ततमो मतः ॥ ६.४७ ॥}
\translate{yōgināmapi sarvēṣāṁ madgatēnāntarātmanā ।\\
śraddhāvānbhajatē yō māṁ sa mē yuktatamō mataḥ ॥ 6.47 ॥}
\cquote{The one with faith, who abides in Me, with their mind focused on Me, and worships Me with unwavering devotion, is considered by Me to be the greatest of all yogis.}

\newpage
\begin{center}
  ॐ तत्सदिति श्रीमद्भगवद्गीतासूपनिषत्सु बह्मविद्यायां योगशास्रे\\
  श्रीकृष्णार्जुनसंवादेध्यानयोगो नाम षष्ठोऽध्यायः॥\\
\arialUC\small ōṃ tatsaditi śrīmadbhagavadgītās ūpaniṣatsu \\brahmavidyāyāṃ yōgaśāstrē\\ śrīkṛṣṇārjunasaṃvādē ātmasaṃyamayōgō \\nāma ṣaṣṭhō'dhyāyaḥ॥
\end{center}
\chapEnd

\newpage
\begin{mananam}{\mananamfont {Mananam sloka - 45, 46}}
\mananamtext What does inner perfection mean to me? How does it differ from outer perfection seen in works, art, or sports?\\
Why is purifying the mind and heart crucial on the spiritual path? How does it impact my growth?\\
How can I improve myself over multiple lifetimes on the spiritual journey? What practices contribute to this evolution?\\
How can I acknowledge my efforts and stay motivated? What mindset helps sustain progress?
\end{mananam}
\WritingHand\enspace{\selfReflect\small{Self Reflection}}
\begin{inspiration}{\mananamfont Inspiration}
\mananamtext The Yoga of self-control and inner purification holds the highest place, according to spiritual teachings. One who has mastered self-control gains immense abilities and can achieve anything effortlessly. A pure heart and a focused mind are invaluable assets for a seeker on the spiritual path. Such spiritual qualities make a person richer than even the materially wealthiest person. 
\end{inspiration}
\newpage